% AA-48: Wymierna gałąź Kovacica dla q w C(x).
% Rozdział samodzielnie dowodzi dokładnej i skończonej procedury
% dla WSZYSTKICH racjonalnych rozwiązań Riccatiego r'+r^2=q.
% Przypadki Kovacica 2 i 3 pozostają wynikami klasycznymi.
\chapter{Wymierna gałąź algorytmu Kovacica dla potencjałów wymiernych}
\chaptermark{Wymierna gałąź Kovacica}
\label{aa:ch:kovacic-rational-potential}

W rozdziale~\ref{aa:ch:risch-kovacic} skonstruowaliśmy
zupełny test wymiernych rozwiązań równania
$r'+r^2=q$ przy założeniu $q\in\mathbb C[x]$.
Dla $q\in\mathbb C(x)$ występują także bieguny
w skończonych punktach, a ich części główne
ograniczają możliwe rozwiązania.
Pokażemy, że również w tym przypadku
\emph{wszystkie wymierne rozwiązania}
można znaleźć przez skończoną liczbę
rozgałęzień algebraicznych i układów
\emph{liniowych}.
Zachowujemy rozróżnienie:
jest to kompletna \emph{pierwsza gałąź}
algorytmu Kovacica, nie pełny algorytm
obejmujący algebraiczne rozwiązania
Riccatiego wyższych stopni.

\section{Konwencja i lokalne dane}

\begin{definicja}[Współczynnik rozwinięcia w nieskończoności]
\label{aa:def:kov-rational-laurent}
Dla $f\in\mathbb C(x)$ przez
$[x^j]_\infty f$ oznaczamy
współczynnik przy $x^j$
w jednoznacznym rozwinięciu
Laurenta funkcji $f$ w nieskończoności.
Dla $j<0$ jest to współczynnik
części właściwej przy
$x^j$, a dla $j\ge0$
współczynnik części wielomianowej.
Przez $\operatorname{pp}_a(f)$
oznaczamy skończoną sumę wyrazów
o ujemnych potęgach $x-a$
w rozwinięciu Laurenta
w skończonym punkcie $a$.
\end{definicja}

\begin{uwaga}[Efektywność obliczeń]
\label{aa:uw:kov-rational-effective}
W twierdzeniu o algorytmie zakładamy,
że współczynniki wielomianów
dopuszczają dokładną arytmetykę,
test równości, rozkład wielomianów
na czynniki nad efektywnym
algebraicznym rozszerzeniem
oraz rozwiązywanie równań
kwadratowych. Wystarcza
na przykład podciało
$\overline{\mathbb Q}$.
Nie utożsamiamy dowolnie
nazwanych liczb zespolonych
bez reprezentacji z danymi
obliczalnymi.
\end{uwaga}

\section{Analiza biegunów skończonych}

\begin{lemat}[Możliwe rzędy biegunów]
\label{aa:lem:kov-rational-local-orders}
Niech $q,r\in\mathbb C(x)$
spełniają $r'+r^2=q$.
W każdym skończonym punkcie
$a\in\mathbb C$ zachodzi:
\begin{enumerate}
\item jeśli $r$ ma biegun rzędu
$m\ge2$, to $q$ ma w $a$
biegun dokładnie rzędu $2m$,
a współczynnik wiodący
$\beta_{-2m}$ jest kwadratem
współczynnika wiodącego
$c_{-m}$ funkcji $r$;
\item jeśli $r$ ma biegun prosty
$r=c_{-1}(x-a)^{-1}+O(1)$,
to $q$ ma biegun rzędu co
najwyżej $2$, a
$[ (x-a)^{-2}]_a q
=c_{-1}(c_{-1}-1)$;
\item jeżeli $q$ jest regularna w $a$,
a $r$ ma tam biegun,
to jest on prosty i ma residuum $1$.
\end{enumerate}
W szczególności, jeśli $q$ ma
w jakimkolwiek skończonym
punkcie biegun nieparzystego
rzędu co najmniej $3$,
równanie nie ma rozwiązania wymiernego.
\end{lemat}
\begin{proof}
Niech $z=x-a$.
Dla $r=c_{-m}z^{-m}+O(z^{-m+1})$,
$m\ge2$, kwadrat $r^2$
ma niezerowy wyraz
$c_{-m}^2z^{-2m}$,
a $r'$ ma biegun rzędu
$m+1<2m$. Wyraz ten
nie może się skasować,
co dowodzi punktu pierwszego.
Dla $m=1$ wyraz przy
$z^{-2}$ w $r'+r^2$
wynosi $-c_{-1}+c_{-1}^2$;
nie ma wyższych ujemnych potęg,
co daje punkt drugi.
Jeżeli $q$ jest regularna,
pierwszy punkt wyklucza
$m\ge2$, a przy $m=1$
otrzymujemy
$c_{-1}(c_{-1}-1)=0$.
Ponieważ $c_{-1}\ne0$,
jest równy $1$.
Jeśli $q$ ma biegun rzędu
nieparzystego co najmniej $3$,
nie może on pochodzić
ani z pierwszego, ani
z drugiego przypadku.
\end{proof}

\begin{lemat}[Skończone listy dopuszczalnych części głównych]
\label{aa:lem:kov-rational-local-candidates}
Niech $a$ będzie biegunem funkcji
$q\in\mathbb C(x)$ stopnia $k$.
Dla każdego wymiernego rozwiązania
$r'+r^2=q$ jego część główna
$\operatorname{pp}_a(r)$ należy
do skończonej, efektywnie
obliczalnej listy:
\begin{enumerate}
\item jeśli $k=1$, lista
ma jeden element $(x-a)^{-1}$;
\item jeśli $k=2$ i
$q=\beta_{-2}(x-a)^{-2}+
O((x-a)^{-1})$, lista składa
się z wyrazów $c(x-a)^{-1}$,
gdzie $c(c-1)=\beta_{-2}$;
\item jeśli $k=2m\ge4$,
lista liczy co najwyżej
dwa wielomiany Laurenta
$\sum_{j=1}^m c_{-j}(x-a)^{-j}$.
Dla każdego z dwóch znaków
$c_{-m}=\pm\sqrt{\beta_{-2m}}$
pozostałe $c_{-m+1},\ldots,c_{-1}$
wyznacza jednoznacznie
porównanie współczynników
przy potęgach
$(x-a)^{-2m+1},\ldots,(x-a)^{-m-1}$
w równaniu
$r'+r^2=q$.
\end{enumerate}
\end{lemat}
\begin{proof}
Jeśli $q$ ma biegun prosty,
to $r$ nie może być regularna,
bo $r'+r^2$ byłaby regularna.
Z Lematu~\ref{aa:lem:kov-rational-local-orders}
wynika prosty biegun $r$,
a brak wyrazu $(x-a)^{-2}$
w $q$ wymusza residuum $1$.
Jeśli $q$ ma biegun drugiego rzędu,
$r$ musi mieć biegun prosty,
którego współczynnik spełnia
podane równanie kwadratowe.

Dla $k=2m\ge4$ biegun $r$
musi mieć rząd $m$.
Niech $u=\operatorname{pp}_a(r)$
i zapiszmy $r=u+v$,
gdzie $v$ jest regularna w $a$.
Wtedy
\[
 r'+r^2=u'+u^2+v'+2uv+v^2.
\]
Ponieważ $u$ ma biegun rzędu $m$,
wyrazy $v'$, $2uv$, $v^2$
mają bieguny rzędu
co najwyżej $m$.
Dlatego współczynniki
przy $z^{-2m},\ldots,z^{-m-1}$
w $q$ zależą wyłącznie
od $u'+u^2$.
Współczynnik przy $z^{-2m}$
daje $c_{-m}^2=\beta_{-2m}$.
W każdym kolejnym równaniu
niewyznaczone dotąd
$c_{-m+j}$ występuje liniowo
ze współczynnikiem
$2c_{-m}\ne0$,
a reszta zależy tylko
od wcześniej wyznaczonych
współczynników.
Tym samym rekurencja
kończy się po $m-1$
krokach i nie wprowadza
dodatkowych wyborów.
\end{proof}

\begin{przyklad}[Biegun czwartego rzędu]
\label{aa:prz:kov-rational-high-pole}
Niech $q=x^{-4}-2x^{-3}$.
W punkcie $0$ mamy $m=2$,
a więc $c_{-2}=\pm1$.
Przy wyborze $c_{-2}=1$
współczynnik przy $x^{-3}$
w $u'+u^2$ wynosi
$-2+2c_{-1}$.
Porównanie z $-2$
daje $c_{-1}=0$,
czyli $u=1/x^2$.
Istotnie,
\[
 (1/x^2)' +(1/x^2)^2
 =-2/x^3+1/x^4=q.
\]
Nie należy traktować
takiego bieguna jako
prostego: jego rząd
determinuje już wiodący
wyraz równania.
\end{przyklad}

\section{Normalna postać wymiernego rozwiązania}

\begin{lemat}[Wydzielenie biegunów ruchomych]
\label{aa:lem:kov-rational-normal-form}
Niech $q,r\in\mathbb C(x)$
i $r'+r^2=q$.
Niech $S\subset\mathbb C$
będzie skończonym zbiorem
wszystkich biegunów $q$.
Dla $a\in S$ połóżmy
$R_a=\operatorname{pp}_a(r)$
oraz $R=\sum_{a\in S}R_a$.
Wówczas istnieją
$\omega\in\mathbb C[x]$
oraz unormowany, niezerowy,
bezkwadratowy $P\in\mathbb C[x]$
takie, że
\begin{equation}\label{aa:eq:kov-rational-normal-r}
 r=\omega+R+\frac{P'}P,
 \qquad
 P(a)\ne0\quad(a\in S).
\end{equation}
Pierwiastki $P$ to dokładnie
bieguny $r$ spoza $S$.
Jeśli $S=\varnothing$, przyjmujemy $R=0$.
\end{lemat}
\begin{proof}
Po odjęciu $R$ funkcja
$r-R$ jest regularna
we wszystkich punktach $S$.
W pozostałych punktach,
w których $q$ jest regularna,
na mocy
Lematu~\ref{aa:lem:kov-rational-local-orders}
każdy biegun $r$ jest prosty
i ma residuum $1$.
Niech $P$ będzie iloczynem
$(x-b)$ po wszystkich takich
pozostałych biegunach $b$;
jeśli ich nie ma, połóżmy $P=1$.
Wielomian $P$ jest unormowany,
bez pierwiastków wielokrotnych,
a żaden jego pierwiastek
nie należy do $S$.
Funkcja $P'/P=\sum_b1/(x-b)$
ma dokładnie części główne
$r-R$ poza $S$.
Zatem $r-R-P'/P$
nie ma skończonych biegunów,
jest więc wielomianem $\omega$.
\end{proof}

\begin{stwierdzenie}[Równanie liniowe dla wielomianu]
\label{aa:st:kov-rational-P-equation}
Niech $q\in\mathbb C(x)$
oraz $u\in\mathbb C(x)$.
Dla każdego niezerowego
$P\in\mathbb C[x]$
równoważne są warunki
\[
 \left(u+\frac{P'}P\right)'
 +\left(u+\frac{P'}P\right)^2=q
\]
oraz
\begin{equation}\label{aa:eq:kov-rational-P-ode}
 P''+2uP'+(u'+u^2-q)P=0.
\end{equation}
Jeśli zarówno $u$, jak i $q$
są regularne w $b$
i $P$ rozwiązuje drugie
równanie, to $P$ nie może
mieć w $b$ pierwiastka
wielokrotnego.
\end{stwierdzenie}
\begin{proof}
Rozwijamy lewą stronę:
\[
 u'+u^2+2u\frac{P'}P+\frac{P''}P;
\]
wyrazy $(P'/P)^2$
z pochodnej ilorazu
i kwadratu znoszą się.
Po pomnożeniu przez $P$
otrzymujemy drugie równanie.
Jeżeli $P$ ma w $b$ zero
rzędu $m\ge2$ i $u$ jest
regularna, to $P''$
ma rząd dokładnie $m-2$,
a pozostałe wyrazy
$2uP'$ i $(u'+u^2-q)P$
mają rzędy co najmniej
$m-1$ i $m$, pod warunkiem,
że $q$ jest regularna
w $b$. W zastosowaniach
poniżej $b$ będzie punktem
spoza biegunów $q$,
więc warunek jest spełniony.
Takie składniki nie mogą się
skasować.
\end{proof}

\section{Kontrola w nieskończoności}

Niech $R$ będzie wybraną sumą
części głównych w punktach
skończonych, a
$\rho=[x^{-1}]_\infty R$.
Wszystkie składniki $R$
są właściwymi funkcjami
wymiernymi, zatem
$R=\rho/x+O(x^{-2})$.
Dla unormowanego $P$
stopnia $n\ge0$ mamy
\begin{equation}\label{aa:eq:kov-rational-logarithmic-infinity}
 \frac{P'}P=\frac n x+O(x^{-2}).
\end{equation}

\begin{lemat}[Skończone możliwości przy nieskończoności]
\label{aa:lem:kov-rational-infinity}
Niech $q\in\mathbb C(x)$,
a $R$ będzie ustaloną sumą
dopuszczalnych części głównych
z Lematu~\ref{aa:lem:kov-rational-local-candidates}.
Załóżmy, że
$r=\omega+R+P'/P$
spełnia $r'+r^2=q$,
gdzie $\omega\in\mathbb C[x]$,
$P$ jest unormowany stopnia $n$
i nie zeruje się w biegunach $q$.
Wtedy zachodzi dokładnie jeden
z poniższych przypadków.
\begin{enumerate}
\item Jeśli część wielomianowa
$q$ ma dodatni stopień,
to jest on parzysty $2d\ge2$,
a $\deg\omega=d$.
Istnieją tylko dwa możliwe
wielomiany $\omega$, wyznaczone
przez warunek
$\deg\operatorname{pp}_\infty(q-\omega^2)
\le d-1$, gdzie
$\operatorname{pp}_\infty$
oznacza część wielomianową,
a ich współczynniki wiodące
są pierwiastkami kwadratowymi
wiodącego współczynnika $q$.
Dla każdego wyboru $\omega$,
gdy $a_d=[x^d]\omega\ne0$,
stopień $P$ jest wymuszony:
\begin{equation}\label{aa:eq:kov-rational-degree-positive}
 n=\frac{[x^{d-1}]_\infty
 \bigl(q-(\omega+R)'-(\omega+R)^2\bigr)}
 {2a_d}.
\end{equation}
\item Jeśli $q=c+O(x^{-1})$
przy $x\to\infty$, $c\ne0$,
to $\omega$ jest stała
i $\omega^2=c$.
Dla każdej z dwóch możliwości
$\omega$ koniecznie
\begin{equation}\label{aa:eq:kov-rational-degree-constant}
 n=\frac{[x^{-1}]_\infty
 \bigl(q-(\omega+R)'-(\omega+R)^2\bigr)}
 {2\omega}.
\end{equation}
\item Jeśli $q=O(x^{-1})$,
to $\omega=0$ i koniecznie
$q=O(x^{-2})$.
Przy $\beta=[x^{-2}]_\infty q$
liczba $n$ musi być
nieujemną liczbą całkowitą
spełniającą
\begin{equation}\label{aa:eq:kov-rational-degree-vanishing}
 (n+\rho)(n+\rho-1)=\beta.
\end{equation}
Istnieją co najwyżej
dwa takie wybory $n$.
\end{enumerate}
\end{lemat}
\begin{proof}
Z~\eqref{aa:eq:kov-rational-logarithmic-infinity}
wynika
$r=\omega+(\rho+n)/x+O(x^{-2})$.
Jeśli $\deg\omega=d\ge1$,
najwyższa potęga w
$r'+r^2$ ma stopień $2d$
i współczynnik będący
kwadratem wiodącego
współczynnika $\omega$.
Wszystkie potęgi
$x^{2d},\ldots,x^d$
pochodzą wyłącznie z
$\omega^2$; człony
$\omega'$, $2\omega R$,
$2\omega P'/P$ mają stopień
co najwyżej $d-1$.
Dopasowujemy więc kolejno
współczynniki $\omega$
od najwyższego w dół,
dla każdego z dwóch
znaków wiodącego pierwiastka.
Przy potędze $x^{d-1}$
pozostały wkład od $P'/P$
wynosi dokładnie $2a_dn$,
a $(P''/P)=O(x^{-2})$.
Stąd~\eqref{aa:eq:kov-rational-degree-positive}.

Jeśli $\omega$ jest niezerową
stałą, to wartość graniczna
$r'+r^2$ wynosi $\omega^2$
i musi być równa $c\ne0$.
Wyraz przy $x^{-1}$ w
równaniu po odjęciu
$(\omega+R)'+(\omega+R)^2$
pochodzi wyłącznie z
$2\omega P'/P$ i jest
równy $2\omega n$.
Daje to drugi wzór.

Jeśli $\omega=0$,
to $r=(\rho+n)/x+O(x^{-2})$.
Zatem $r'+r^2=O(x^{-2})$,
a współczynnik przy $x^{-2}$
wynosi
$-(\rho+n)+(\rho+n)^2$.
To jest dokładnie
\eqref{aa:eq:kov-rational-degree-vanishing}.
Równanie kwadratowe
ma najwyżej dwa rozwiązania.
Wymienione możliwości
wyczerpują wszystkie
stopnie części wielomianowej $r$.
\end{proof}

\section{Główne twierdzenie algorytmiczne}

\begin{twierdzenie}[Zupełny test wymiernej gałęzi Kovacica]
\label{aa:tw:kov-rational-complete-algorithm}
Niech $q\in\mathbb C(x)$,
a współczynniki będą podane
efektywnie w sensie
Uwagi~\ref{aa:uw:kov-rational-effective}.
Istnieje skończony algorytm
rozstrzygający, czy
\begin{equation}\label{aa:eq:kov-rational-main-riccati}
 r'+r^2=q
\end{equation}
ma rozwiązanie $r\in\mathbb C(x)$.
Jeżeli tak, algorytm
wyznacza co najmniej jedno
takie rozwiązanie, a następnie
lokalne niezerowe rozwiązanie
$z''=qz$ postaci
$z=\exp(\int r\,\mathrm dx)$.
\end{twierdzenie}
\begin{proof}
\emph{Krok 1: bieguny skończone.}
Rozkładamy mianownik $q$
i znajdujemy jego skończony
zbiór biegunów $S$
oraz ich rzędy.
Jeśli którykolwiek ma
nieparzysty rząd co najmniej
$3$, zgłaszamy
brak rozwiązania wymiernego,
zgodnie z
Lematem~\ref{aa:lem:kov-rational-local-orders}.
W przeciwnym razie dla
każdego $a\in S$ wyznaczamy
jedną lub dwie możliwe
części główne z
Lematu~\ref{aa:lem:kov-rational-local-candidates}.
Wykonujemy wszystkie skończenie
wiele kombinacji tych wyborów,
a dla każdej z nich
definiujemy $R=\sum_aR_a$.

\emph{Krok 2: nieskończoność.}
Dla wybranego $R$
stosujemy
Lemat~\ref{aa:lem:kov-rational-infinity}.
Uzyskujemy skończoną listę
par $(\omega,n)$.
Odrzucamy wszystkie,
dla których wymagany
stopień $n$ nie jest
nieujemną liczbą całkowitą.
Warunki z lematu są konieczne,
więc niczego nie tracimy.

\emph{Krok 3: liniowy układ dla $P$.}
Dla każdej pozostałej pary
stawiamy $u=\omega+R$ i
$P=x^n+\sum_{j=0}^{n-1}p_jx^j$.
Równanie~\eqref{aa:eq:kov-rational-P-ode},
po pomnożeniu przez
wspólny mianownik,
daje skończenie wiele równań
\emph{liniowych} na
$p_0,\ldots,p_{n-1}$.
Rozstrzygamy ich zgodność
dokładną eliminacją liniową.
Nie wystarcza wybrać
dowolnego rozwiązania
układu: dodatkowo żądamy
$P(a)\ne0$ dla każdego
$a\in S$.

Warunki te także
można sprawdzić efektywnie:
rozwiązania zgodnego
układu liniowego tworzą
afiniczną przestrzeń $V$.
Dla każdego $a\in S$
równanie $P(a)=0$
wyznacza podzbiór afiniczny
będący hiperpłaszczyzną
lub całym $V$.
Jeżeli dla pewnego $a$
całe $V$ jest w niej zawarte,
odrzucamy daną parę
$(\omega,n)$.
Jeśli dla żadnego $a$
tak nie jest, skończona
suma właściwych
hiperpłaszczyzn nie pokrywa
przestrzeni afinicznej
nad nieskończonym ciałem
$\mathbb C$. Możemy więc
wybrać $P\in V$ bez zer
w $S$, na przykład
sprawdzając kolejno
skończenie wiele
specjalizacji parametrów
w efektywnym nieskończonym
podciele współczynników.
Wszystkie pierwiastki $P$
poza $S$ są automatycznie
proste: gdyby $b\notin S$
był pierwiastkiem krotności
co najmniej $2$,
w równaniu~\eqref{aa:eq:kov-rational-P-ode}
człon $P''$ miałby
niekasowalny wyraz wiodący
rzędu o dwa mniejszego
niż rząd zera $P$,
bo $u$ i $q$ są w $b$
regularne.

\emph{Poprawność wyniku pozytywnego.}
Jeśli znaleziono $P$,
to $r=u+P'/P$ jest
funkcją wymierną i
ze Stwierdzenia
\ref{aa:st:kov-rational-P-equation}
spełnia dokładnie
\eqref{aa:eq:kov-rational-main-riccati}.
Z Lematu~\ref{aa:lem:kovacic-riccati}
dostajemy lokalne niezerowe
rozwiązanie $z''=qz$
po dołączeniu elementu
$z$ o $z'/z=r$.

\emph{Zupełność wyniku negatywnego.}
Niech $r\in\mathbb C(x)$
będzie dowolnym rozwiązaniem.
Na mocy
Lematów~\ref{aa:lem:kov-rational-local-orders}
i~\ref{aa:lem:kov-rational-local-candidates}
jego części główne pojawiają
się w jednej z badanych
kombinacji $R$.
Z Lematu~\ref{aa:lem:kov-rational-normal-form}
otrzymujemy $r=\omega+R+P'/P$
z unormowanym $P$
bez zer na $S$.
Lemat~\ref{aa:lem:kov-rational-infinity}
zapewnia, że odpowiednie
$(\omega,\deg P)$ jest
na naszej skończonej liście.
Wreszcie $P$ spełnia
liniowe równanie
\eqref{aa:eq:kov-rational-P-ode}
oraz wszystkie warunki
$P(a)\ne0$.
Algorytm zatem znajdzie
co najmniej jedną
poprawną parę i nie
może zakończyć się odpowiedzią
negatywną.
Dowiedliśmy obu kierunków
oraz skończoności całej procedury.
\end{proof}

\begin{uwaga}[Zakres wyniku]
\label{aa:uw:kov-rational-not-all}
Twierdzenie
\ref{aa:tw:kov-rational-complete-algorithm}
rozstrzyga istnienie
\emph{wymiernej} logarytmicznej
pochodnej rozwiązania
$z''=qz$ dla każdego
$q\in\mathbb C(x)$.
Nie rozstrzyga samo
istnienia rozwiązania
liouvillowskiego w ogólności:
gałęzie algorytmu Kovacica,
w których pochodna
logarytmiczna jest algebraiczna
stopnia $2$, $4$, $6$ lub $12$,
wymagają dalszych procedur.
Ogólną klasyfikację przytoczono
w~\ref{aa:tw:kovacic-classification}
jako twierdzenie klasyczne.
\end{uwaga}

\section{Przykłady z pełnymi rachunkami}

\begin{przyklad}[Biegun prosty potencjału]
\label{aa:prz:kov-rational-simple}
Weźmy
\[
 q=x^2+3+\frac2{x-1}.
\]
Dla $a=1$ potencjał ma
biegun prosty, więc
$R=1/(x-1)$.
Jego część wielomianowa
ma stopień $2$, zatem
$\omega=\pm x$.
Dla $\omega=x$ mamy
$u=x+1/(x-1)$ oraz
\[
 u'+u^2
 =x^2+3+\frac2{x-1}=q.
\]
Lemat o nieskończoności
daje $n=0$ i $P=1$.
Zatem $r=x+1/(x-1)$,
a lokalnym rozwiązaniem
jest $z=(x-1)e^{x^2/2}$.
Sprawdzenie $z''/z=q$
można wykonać bezpośrednio
lub przez równanie Riccatiego.
\end{przyklad}

\begin{przyklad}[Dwa możliwe residua]
\label{aa:prz:kov-rational-double}
Niech $q=2/x^2$.
Przy $a=0$ jest biegun
rzędu $2$ z
$\beta_{-2}=2$.
Równanie $c(c-1)=2$
daje $c=2$ albo $c=-1$.
Dla $R=2/x$ przy
nieskończoności mamy
$\rho=2$, $\beta=2$;
równanie
$(n+2)(n+1)=2$
ma dopuszczalne $n=0$.
Z $P=1$ otrzymujemy
$r=2/x$.
Analogicznie drugi wybór
$R=-1/x$, $n=0$
daje $r=-1/x$.
Oba spełniają
\[
 r'+r^2=\frac2{x^2}.
\]
Odpowiadają rozwiązaniom
$z_1=x^2$ oraz $z_2=x^{-1}$
równania $z''=(2/x^2)z$.
\end{przyklad}

\begin{przyklad}[Brak rozwiązania i przypadek algebraiczny]
\label{aa:prz:kov-rational-negative}
Dla $q=1/x^3$
biegun w $0$ ma rząd
nieparzysty większy od $2$,
więc żadne rozwiązanie
Riccatiego nie jest wymierne.
Nie wynika z tego
nieistnienie wszystkich
rozwiązań liouvillowskich.

Dla $q=x+5/(16x^2)$
pierwsza gałąź również
jest wykluczona,
ponieważ część wielomianowa
$q$ ma nieparzysty
dodatni stopień.
Jednocześnie
$r=\sqrt{x}-1/(4x)$
jest algebraicznym,
niewymiernym rozwiązaniem
Riccatiego i prowadzi
do liouvillowskiego
rozwiązania $z$ opisanego
w Przykładzie
\ref{aa:prz:kovacic-quadratic-example}.
Pokazuje to konieczność
rozróżniania odpowiedzi
pierwszej gałęzi
od odpowiedzi całego
algorytmu Kovacica.
\end{przyklad}

\begin{przyklad}[Biegun czwartego rzędu i zachowanie w nieskończoności]
\label{aa:prz:kov-rational-fourth}
Dla $q=x^{-4}-2x^{-3}$
wybór $R=1/x^2$
pochodzi z analizy
bieguna czwartego rzędu.
Mamy $\rho=0$, a
$q=O(x^{-3})$, więc
$\beta=[x^{-2}]_\infty q=0$.
Równanie
$n(n-1)=0$ daje
kandydatów $n=0,1$.
Przy $n=0$ można
wybrać $P=1$,
a wówczas $r=1/x^2$
jest rozwiązaniem.
Wystarcza jeden
pozytywny kandydat;
algorytm nie musi
wyliczać wszystkich
możliwych rozwiązań,
aby rozstrzygnąć istnienie.
\end{przyklad}

\section{Kontrola połączenia z teorią Galois}

Jeżeli algorytm zwraca
wymierne $r\in\mathbb C(x)$,
to w rozszerzeniu
Picarda--Vessiota równania
$z''=qz$ istnieje
niezerowe rozwiązanie
o logarytmicznej pochodnej
$r$. Jego jednowymiarowa
przestrzeń rozwiązań
jest zachowywana przez
grupę Galois, więc
reprezentacja grupy
w odpowiedniej bazie
jest trójkątna.
Jest to kierunek
\emph{od wymiernego Riccatiego
do redukowalności reprezentacji}.
Przejście odwrotne
w ogólnym języku teorii
Picarda--Vessiota korzysta
z klasycznego twierdzenia
o niezmienniczej prostej
i nie jest tu nowym
dowodem strukturalnym
pełnej teorii Galois.

Dalsze przypadki
imprymitywne i skończone
należą do pełnego algorytmu
Kovacica, a nie do
niniejszego rozdziału.
